
\section{Related Work}
\label{sec:related}

HALO sits at the intersection of four literatures, and in each of them some part of what we do
already exists. We name the closest work in each and state what is left.

\textbf{Cross-dataset and foundation models for IMU sensing.}
Deep HAR has moved from single-dataset CNN--LSTM models~\cite{deepconvlstm} through methods that
attack one axis of heterogeneity at a time---contrastive fusion for small-data multimodal
HAR~\cite{cosmo}, physics-informed augmentation for cross-user and cross-device
transfer~\cite{unihar}, federated modelling of inter-user structure~\cite{clusterfl}---to
self-supervised encoders intended to transfer: LiMU-BERT~\cite{limubert} pretrains by span-masked
reconstruction at $20$\,Hz, CrossHAR~\cite{crosshar} combines masked reconstruction with
contrastive learning, and ssl-wearables~\cite{sslwearables} scales multi-task self-supervision to
roughly $700$k person-days of wrist accelerometry. A second wave targets foundation-scale
wearable modelling: LSM~\cite{lsm1} establishes the label-efficiency argument (up to $8\times$
fewer labelled examples), LSM-2~\cite{lsm2} learns directly from incomplete recordings,
RelCon~\cite{relcon} learns a motion-similarity metric over $\sim\!10^9$ windows, and
SensorLM~\cite{sensorlm} pairs sensor data with hierarchical captions and reports zero-shot,
few-shot, and cross-modal retrieval together. General time-series models~\cite{moment,patchtst}
and tabular ones~\cite{mitra} trade domain-specific bias for breadth. A recent survey maps this
space~\cite{harfmsurvey}, and Inertia-1~\cite{inertia1} runs a controlled lifecycle study over
modality, placement, rate, and window length. \emph{What is left:} all of these treat sampling
rate as a preprocessing decision to be normalised away. Our argument is that this is where the
information loss happens, and Section~\ref{sec:bias_experiments} measures it rather than
asserting it.

\textbf{Sampling-rate and resolution invariance.}
This is the literature closest to our L1 mechanism, and it is not in HAR. Rate-invariant
autoencoding~\cite{rateinvae} learns representations quotiented by temporal warping;
FlowState~\cite{flowstate} builds a sampling-rate-\emph{equivariant} forecasting model on a
state-space encoder with a functional-basis decoder; universal spectral
tokenization~\cite{spectraltok} ingests native, irregular wavelength grids without resampling and
explicitly proposes the same idea for multiresolution time series. Within IMU sensing,
MoPFormer~\cite{mopformer} tokenizes into learned motion primitives to make datasets comparable,
but resamples every dataset to $100$\,Hz before quantising, so its codebook is defined over a
fixed-rate grid. \emph{What is left:} FlowState conditions on the rate---it is equivariant,
adapting its computation---whereas HALO's features are \emph{anchored} to physical Hz, so a band
denotes the same physical frequency for every device and no conditioning is required to compare
across them; and the astronomical setting of~\cite{spectraltok} has no gravity, no Nyquist
guard tied to a wearable's firmware, and no anti-aliasing consequence to demonstrate. To our
knowledge no prior HAR system claims rate-invariance by construction, and none measures what
resampling costs on held-out data---which is the experiment we run, and which we regard as the
load-bearing evidence for L1 rather than the mechanism itself.

\textbf{Language as an interface to sensors.}
Two distinct uses of language must be separated. On the \emph{output} side, aligning a modality
to sentence embeddings for open-vocabulary recognition descends from
DeViSE~\cite{devise}, ConSE~\cite{conse}, and CLIP~\cite{clip}; ImageBind~\cite{imagebind} extends
it to IMU among other modalities, and LanHAR~\cite{lanhar} and LLaSA~\cite{llasa} bring it to
HAR. On the \emph{input} side---describing the sensor rather than the activity---the lane is
crowded and recently so. GOAT~\cite{goat} supervises with natural language over both activity
labels and device on-body positions. ZeroHAR~\cite{zerohar} conditions on textual sensor
context---sensor type, Cartesian axis, body position---and aligns the motion latent to it before
aligning to activity text. CHARM and its successor~\cite{charm,sensorsvoice} feed channel-level
textual descriptions into a channel-order-equivariant transformer, and their ablation finds that
conditioning on channel identity via text is one of the largest performance drivers.
SensorLLM~\cite{sensorllm} generates per-channel descriptions to teach an LLM the sensor layout;
oneHAR~\cite{onehar} handles variable combinations of sensor positions with LLM assistance;
UniMTS~\cite{unimts} takes the complementary route of encoding placement structurally as a joint
graph with rotation-invariant augmentation. \emph{What is left, honestly:} ``sensor identity as
natural language'' is not ours to claim. Our differences are specific and narrower than the
general idea. First, granularity: ZeroHAR conditions on a fixed, closed set of categorical
attributes, whereas ours is free text and therefore composes to descriptions never seen; the
channel-description line~\cite{charm,sensorsvoice,sensorllm} is free text but attaches it
per-channel, which is exactly the configuration whose corpus-identifiability we argue against and
test in Table~\ref{tab:identity_probe}. Second, what the conditioning feeds: in all of the above
it feeds a parametric head, whereas ours feeds a non-parametric memory. Third, invariance versus
conditioning: UniMTS's rotation invariance discards orientation, and with it gravity and posture,
which conditioning preserves.

\textbf{Retrieval, memory, and few-shot recognition.}
Deciding by comparison to stored examples is old in this field and new in its current form. The
ubicomp lineage recognises gestures by template: uWave~\cite{uwave} needs a single example per
gesture and matches by dynamic time warping; Junker et al.~\cite{junker} preselect candidate
segments by similarity search before classifying, which is structurally the retrieve-then-decide
shape we use; WarpingLCSS~\cite{warpinglcss} does robust online template spotting from
crowdsourced annotations. The modern lineage is episodic and parametric---matching
networks~\cite{matchingnet}, prototypical networks~\cite{protonet}---and, in language and vision,
non-parametric memory attached to a frozen model~\cite{knnlm,rag,retro,knnvision,reco}.
Memory-modular classification~\cite{memorymodular} states our Phase-B property as its
contribution: generalise to new classes by replacing memory contents, without retraining. In
sensing specifically, RAG-HAR~\cite{raghar} performs training-free retrieval over signal
descriptors and lets an LLM decide, recognising unseen activities; ZARA~\cite{zara} retrieves
multi-sensor evidence into an LLM agent, under the same ``evidence-grounded'' name we use for a
different mechanism; and in-context time-series classifiers~\cite{timee,ticfm} take a labelled
support set and a query in one forward pass with no gradient step, collapsing zero-shot and
few-shot into a single inference mode. \emph{What is left:} not ``we retrieve.'' The specific
claim is that retrieval is performed in a space where similarity is \emph{sensor} similarity
between encoded exemplar \emph{patches}, produced by an encoder that was conditioned on
acquisition configuration---which is what lets a wrist query at an unusual rate retrieve
comparable evidence, and which descriptor-space~\cite{raghar}, LLM-agent~\cite{zara}, and
CLIP-space~\cite{sensorlm} retrieval cannot do. To that we add calibrated
rejection~\cite{selectiveclass,edl}, which none of the retrieval-based HAR systems above
provide. Section~\ref{sec:label_efficiency} reports the control that decides whether this
distinction is real: the same memory mechanism on a competitor's frozen encoder.

\textbf{Self-supervised objectives.}
Phase~A composes established pieces: contrastive~\cite{simclr,infonce} and supervised
contrastive~\cite{supcon} learning as controls; non-contrastive objectives that avoid declaring
in-batch negatives~\cite{vicreg,barlowtwins}; masked modelling from language and
vision~\cite{bert,mae}; and latent-target prediction against a momentum
teacher~\cite{byol,data2vec,ijepa}. Our time--frequency term is inspired by TF-C~\cite{tfc} but
is not faithful to it, and we say so where it is described. The contribution here is not a new
objective but the choice to remove activity labels from all of them, and the cross-placement
positives that a corpus with verified simultaneous multi-device recordings makes available.

\subsection{The Combination, and Why It Is Not a List}
\label{sec:combination}

Table~\ref{tab:coverage} summarises which axes each of the closest systems explicitly addresses.
No prior system addresses all of them, which is true but is also the least interesting thing to
say about it: ``we combine more components'' is not a contribution, and a reviewer is right to
be unmoved by a filled-in row. Variable channel counts have known solutions; variable sampling
rates have known solutions; language as a bridge to an open label vocabulary has known
solutions. We claim none of those as novel.

We make a different kind of claim: that physical-Hz \emph{anchoring} is what makes a single,
flat, cross-device memory well-posed in the first place---an \emph{enabling} result rather than
an assembly. A non-parametric memory shared across deployments requires that a query encoded on
one device and an exemplar encoded on another be comparable \emph{as vectors}, with no
per-configuration index and no conditioning step to reconcile them. Resampling does not provide
this (it discards content the two devices do not share), and rate-equivariance does not provide
it either (it relates the two by a learned transformation rather than making them the same
quantity). This is why the existing retrieval-based sensing systems retrieve somewhere
else---in hand-designed descriptor space~\cite{raghar}, through an LLM
agent~\cite{zara}, or in CLIP space~\cite{sensorlm}---and why none of them needs the property
we are claiming. Section~\ref{sec:bias_experiments} tests it directly, using recordings in
which two devices captured the same person performing the same activity at the same instant, so
that configuration is the only variable (Table~\ref{tab:crossconfig}).

Once that claim is in view, the remaining choices stop looking like options. They are
\emph{mutually entailed}---each one is what makes another valid---so the set is a consequence of
L1 and L2 rather than a collection of independently attractive features. And because ``it is
just a combination'' is an empirical objection, we answer it empirically as well: the
conventional assembly of the same ingredients---resample to a common rate, attach free text per
channel, put retrieval on top---is run as a rival arm in
Table~\ref{tab:tokenizer_ablation}, not merely argued against.

\begin{table}[t]
  \centering
  \resizebox{\columnwidth}{!}{
  \footnotesize
  \setlength{\tabcolsep}{4pt}
  \begin{tabular}{l ccccc}
    \toprule
    & \multicolumn{2}{c}{acquisition} & orient./ & label &$k$-curve \\
    \cmidrule(lr){2-3}
    System & rate & placement & gravity & interface & (no grad.) \\
    \midrule
    LiMU-BERT~\cite{limubert}, CrossHAR~\cite{crosshar} &        &        &        &        &  \\
    harnet~\cite{sslwearables}, LSM~\cite{lsm1,lsm2}    &        &        &        &        & \checkmark \\
    NormWear~\cite{normwear}                            &        & $\circ$ &    & \checkmark &  \\
    UniMTS~\cite{unimts}                                &        & \checkmark & $\circ$ & \checkmark &  \\
    ZeroHAR~\cite{zerohar}                              &        & $\circ$ &        & \checkmark &  \\
    GOAT~\cite{goat}, oneHAR~\cite{onehar}              &        & \checkmark &        & \checkmark &  \\
    CHARM~\cite{charm,sensorsvoice}                     &        & \checkmark &        &        &  \\
    MoPFormer~\cite{mopformer}                          &        & $\circ$ &        &        &  \\
    FlowState~\cite{flowstate}                          & $\circ$ &    &        &        &  \\
    RAG-HAR~\cite{raghar}, ZARA~\cite{zara}             &        &        &        & \checkmark & \checkmark \\
    SensorLM~\cite{sensorlm}                            &        &        &        & \checkmark & \checkmark \\
    \midrule
    \textbf{HALO}                                       & \checkmark & \checkmark & \checkmark & \checkmark & \checkmark \\
    \bottomrule
  \end{tabular}}
  \caption{Which axes each system \emph{explicitly claims} to address. $\checkmark$~= addressed
  in the form we argue is needed; $\circ$~= addressed in a weaker or lossy form (rate
  \emph{equivariance} rather than physical-Hz anchoring; \emph{categorical} rather than free-text
  placement; rotation \emph{invariance}, which discards gravity and therefore posture); blank =
  not claimed by that work. A blank is not a criticism---these are different papers solving
  different problems---and the table reports stated scope, not measured performance. The point
  last column marks works that report performance across a range of label budgets; HALO's mark
  there differs in kind rather than degree, in that one frozen model traverses every budget with
  no gradient update at any point, whereas the other marked rows fit something at each budget.
  The point of the table is not the bottom row's length; it is the dependency structure
  described below.}
  \label{tab:coverage}
  \vspace{-1em}
\end{table}

\textbf{Rate anchoring is a precondition for the memory, not a parallel feature.} A
non-parametric memory only works if a query and a stored exemplar recorded on different devices
are comparable \emph{as vectors}. If the encoder is rate-\emph{equivariant}---told the rate and
adapting---then two windows at different rates are related by a learned transformation, and the
memory must either be indexed per configuration or trust that transformation to be exact. Because
HALO's band $k$ denotes the same physical frequency for every device, one flat memory is
well-posed. L1 exists in the form it does because of Phase~B.

\textbf{The granularity of the language conditioning is forced by the retrieval, too.}
Per-channel free text~\cite{charm,sensorsvoice,sensorllm} is the natural choice if the
conditioning feeds a classifier: it is maximally expressive and the head can ignore what it does
not need. It is the wrong choice if the conditioning feeds a memory, because a description that
is nearly unique to a corpus makes \emph{same-corpus} neighbours the nearest ones, which is
precisely the shortcut cross-dataset evaluation exists to catch. Attaching identity to the
\emph{sensor} keeps the description expressive enough to name an unseen device and too coarse to
name a dataset. That is why Table~\ref{tab:identity_probe} is in the paper.

\textbf{Preserving gravity is forced by the geometry finding.} Rotation
invariance~\cite{unimts} is the standard answer to unknown mounting, and it is lossy in one
specific way: it discards the gravity direction, and with it posture. Section~\ref{sec:geometry}
shows that the label pairs a text projection cannot separate are exactly the posture-versus-
transition pairs (\emph{sitting}/\emph{sitting down}, \emph{lying}/\emph{standing}). A system
that gives up on the text geometry \emph{must} therefore keep posture in the signal, so
conditioning rather than invariance is not a preference here but a requirement.

\textbf{Label-free pretraining is forced by the same shortcut argument.} Any closed-vocabulary
objective over a multi-corpus training set rewards inferring the corpus and predicting its
marginal. That is survivable for a parametric head, which can be refitted per deployment; it is
not survivable for a frozen encoder feeding a memory, because the shortcut is baked into the
stored vectors.

So the four choices form a chain: removing any one of them invalidates another. That is a
different and more falsifiable claim than ``we do all of these,'' and
Section~\ref{sec:bias_experiments} is built to test it---each ablation in
Table~\ref{tab:tokenizer_ablation} removes one link and asks whether the rest still stand.

\textbf{A note on scale.} We do not think the scale-first programme is naive, and we do not
claim it fails on its own terms. On population-scale health outcomes, models trained on tens of
millions of hours~\cite{sslwearables,lsm1,sensorlm} are the right instrument and we would not
beat them. The disagreement is about which axis a corpus can close. Scale reduces surprise along
the axes a corpus \emph{samples}---subjects, activities, sessions. It cannot close the axes that
are \emph{generated at deployment time} by a continuous, composable chain of mounting,
placement, filtering, and convention choices, because those are not drawn from the corpus
distribution at all. Our position is not ``less data, cleverer model''; it is that these two axes
call for different remedies, and that the second one has been under-served because it is
invisible to the benchmark protocols the field uses.
